\Titular*%
{filosofia}%
{Podemos chegar á verdade científica?}%
{Mauro Garrido Rodríguez}%
{Onde a física e a metafísica se confunden.}%

\begin{multicols}{2}

\subsection*{Que é a verdade? As verdades científicas ou \textit{<<da
existencia>>}}

É posible realmente coñecer a verdade? Podemos saber que algo será verdade
sempre? Dito doutro xeito: existe a verdade, e de existir, poden os seres
humanos coñecela?

Reflexionemos primeiro que entendemos por verdade. Parece claro que unha
verdade, para ser tal, debe selo sempre, baixo calquera circunstancia e para
todo o mundo: as verdades circunstanciais non son universais, e polo tanto hai
algo \textit{verdadeiro} que se lles escapa, polo que non son verdades en
rigor. Así, a verdade fala da \textit{esencia}: do que as cousas son en si; non
do que aparentan, non do que nos mostran, senón do que son.

Inspirándonos no filósofo Karl Jaspers \cite{karl}, podemos facer unha
distinción xeral entre dous tipos de verdade: as do existente e as da
existencia.

As da existencia ou <<científicas>> son as que trataremos neste artigo: sono
sobre o mundo, sobre o universo; analizamos o mundo natural e, en base á
experiencia, enunciamos postulados, leis e conclusións derivadas destes. Así,
parte da verdade científica é o que a experiencia demostra estatisticamente, na
repetición dos nosos experimentos. Verdade é dicir: os corpos caen cara á
superficie da Terra, a luz crea un padrón de difracción nunha pantalla ao pasar
por unha fenda, o Sol ponse polo horizonte, ... Porque en todos estes casos se
repetimos o experimento, volveremos obter o mesmo resultado.

\begin{figure}[H]
    \centering
    \includegraphics[width=\linewidth]{revistas/003/imaxes/atardecer2.jpg}
    \caption{Atardecer sobre as Illas Cíes, Vigo (Fonte: imaxe persoal)}
\end{figure}

\subsection*{A verdade ontolóxica}

Pero a física vai máis alá. Busca, non só describir o mundo natural, senón
tamén predicilo (construír modelos que funcionen), e, non contenta con isto,
\textbf{descubrir que o constitúe e como é}, a un nivel fundamental.
Recuperando os exemplos que puxemos antes: a forza da gravidade é a responsable
da atracción entre corpos con masa, a luz é unha onda, a Terra é un xeoide e
xira arredor de si mesma cunha determinada inclinación con respecto ao plano da
eclíptica mentres orbita arredor do Sol... Agora non só describimos o que
sucede, senón que damos unha imaxe do mundo e dos elementos que o compoñen, de
xeito que baixo estes supostos \textit{ontolóxicos} (do grego antigo
\textit{onto-logos}, relativos ao estudo do \textit{ente}, do que
\textit{<<é>>}) podemos describir e predicir o que sucederá. Isto é, a verdade,
o real para a ciencia.

\begin{figure}[H]
    \centering
    \includegraphics[width=\linewidth]{revistas/003/imaxes/difraccion2.jpg}
    \caption{
        Difracción da luz no laboratorio de Técnicas Experimentais III (Fonte:
        imaxe persoal)
    }
    \label{fig:placeholder}
\end{figure}

\subsection*{As verdades científicas tamén inclúen presupostos
\textit{<<metafísicos>>}}

É máis: detrás destes tres exemplos hai un presuposto que tradicionalmente
chamariamos <<filosófico>>, que non demostramos; e que, no entanto, empregamos
non só nestes exemplos, senón durante toda a nosa historia ata o século XX: o
principio metafísico de \textbf{\textit{contrafactualidade}} \cite{bernard}:
asumimos que tódalas variables tiñan valores definidos aínda cando non as
mediamos, que o campo electromagnético estaba en tódolos puntos do espazo... En
xeral, facemos un suposto contrafactual cando dicimos que \textit{<<un sistema
físico tería propiedades ben definidas aínda cando non esteamos medindo as
devanditas propiedades>>}.

Coa chegada da cuántica, este suposto, que en principio parecía tan evidente
que se admitía de xeito implícito, quedou contra as cordas, aínda que na
macroescala se cumpre en boa aproximación. A nivel atómico as partículas non
seguen traxectorias definidas e, por exemplo, non é posible asignar unha única
velocidade para unha partícula nun mesmo estado e nun mesmo experimento: antes
de facer a medición, non tiña un valor definido da súa velocidade.

As desigualdades de Bell achegan aínda máis evidencia: unha situación na cal,
nun sistema de partículas entrelazadas, asumir uns valores definidos (as
famosas \textit{variables ocultas}) leva a resultados que non se corresponden
co que se obtén experimentalmente (se ben é certo que estas teorías de
variables ocultas, que son contrafactuais, poden rescatarse se admitimos
influencias máis rápidas que a luz, aínda que quedan en dúbida).

\begin{figure}[H]
    \centering
    \includegraphics[width=\linewidth]{revistas/003/imaxes/popper.jpg}
    \caption{Karl Popper, filósofo da ciencia (Fonte: \cite{karl_pop})}
    \label{popper}
\end{figure}

Así, a física tamén depende de \textit{<<verdades metafísicas>>}, como o
presuposto fundamental necesario para a mesma idea de ciencia: que existe un
mundo real ata certo punto independente do observador (realismo ontolóxico),
que este mundo natural ten regularidades estables e que é \textit{intelixible},
que a mente humana pode coñecer estas regularidades.

\subsection*{Ciencia, filosofía e falsabilidade}

Por suposto, tamén temos que ser conscientes das limitacións da ciencia: cando
os nosos postulados predín ben e non fallan ante ningún experimento ,dicimos que
chegamos á verdade, e que a nosa imaxe do mundo é representativa de como
funciona a realidade. Mais isto é, por suposto, unha ilusión: xamais poderemos
estar seguros de ter a última palabra; sempre pode haber un experimento que
negue as nosas premisas ou que evidencie fenómenos aínda non explicados. Así, a
ciencia persegue a verdade, mais sen nunca chegar a atrapala. Por iso, neste
artigo sempre falamos de \textit{<<falsabilidade>>} e non de
\textit{<<verificabilidade>>}, porque xamais é posible \textit{verificar} un
determinado enunciado científico, como ben nos fixo ver Karl Popper.

Mais isto da falsabilidade, \textbf{e que a ciencia só se debe ocupar do
falsable}, é moitas veces utilizado para menosprezar o papel da filosofía na
ciencia... Sendo certo, non é argumento para desdeñala, posto que ademais de
ter a capacidade de guiar a empresa científica, non remata aí o seu papel: en
moitas ocasións un presuposto metafísico que xulgabamos imposible de pór a
proba acaba resultando ser falsable, ou polo menos posto en dúbida pola
evidencia.

A existencia dos átomos tamén se consideraba un presuposto indemostrable ata
hai relativamente pouco, cando físicos e filósofos como Mach e Ostwald
sinalaban o atomismo como unha \textit{metafísica inútil}.

\nocite{mermin}

\printbibliography

\end{multicols}
